\clearpage\SourceLeaf{85}%
\pdfbookmark[1]{Berechnung des jüdischen Osterfestes}{osterfest1802}
\begin{center}
{\large BERECHNUNG DES JÜDISCHEN OSTERFESTES.\par}
\medskip\hrule\medskip
{\small Monatliche Correspondenz zur Beförderung der Erd- und Himmels-Kunde,\par
herausgegeben vom Freiherrn \textsc{von Zach}. Mai 1802.\par}
\medskip\hrule
\end{center}\medskip
Der 15. Nisan des jüdischen Jahrs $A$, an welchem die Juden ihr Osterfest
feiern, fällt in das Jahr $A-3760=B$ der christlichen Zeitrechnung; zur Bestimmung
des entsprechenden Monatstages dient folgende rein arithmetische
Regel:

Man dividire $12A+17$, oder welches hier einerlei ist, $12B+12$ mit 19,
und nenne den Rest $a$; ferner dividire man $A$ oder $B$ mit 4, und setze den
Rest $=b$. Man berechne den Werth
\begin{center}\small
\begin{tabular}{lll}
\multicolumn{1}{c}{von} & \multicolumn{1}{c}{oder von} & \shortstack{Werth der Decimalbrüche\\in gemeinen Brüchen}\\[5pt]
$32.0440932$ & $20.0955877$ & $32\frac{4343}{98496}\quad20\frac{9415}{98496}$\\[5pt]
$+1.5542418a$ & $+1.5542418a$ & $1\frac{272953}{492480}$\\[5pt]
$+0.25\quad b$ & $+0.25\quad b$ & $\frac14$\\[5pt]
$-0.003177794A$ & $-0.003177794B$ & $\frac{313}{98496}$
\end{tabular}
\end{center}
und setze ihn $=M+m$, so dass $M$ die ganze Zahl und $m$ den (Decimal-)
Bruch bedeute. Endlich dividire man $M+3A+5b+5$ oder $M+3B+5b+1$
mit 7, und setze den Rest $=c$. Nun hat man folgende vier Fälle zu unterscheiden:

I. Ist $c=2$ oder 4 oder 6, so fällt Ostern den $M+1^{\mathrm{ten}}$ März \textit{alten Styls},
wofür man den $M-30^{\mathrm{ten}}$ April schreibt, wenn $M>30$ wegen \textit{Adu}.
\SourceLeaf{86}%

II. Ist $c=1$, zugleich $a>6$ und ausserdem $m\geqq0.63287037$
$(\frac{311676}{492480}=\frac{1367}{2160})$,
so fällt Ostern den $M+2^{\mathrm{ten}}$ März \textit{a. St.} wegen \textit{Getred}.

III. Ist $c=0$, zugleich $a>11$ und noch $m\geqq0.89772376$
$(\frac{442111}{492480}=\frac{23269}{25920})$,
so ist Ostern den $M+1^{\mathrm{ten}}$ März \textit{a. St.} wegen \textit{Batu Thakpad}.

IV. In allen übrigen Fällen ist Ostern den $M^{\mathrm{ten}}$ März alten Styls.

\begin{center}\rule{.25\linewidth}{.3pt}\end{center}

\textit{Erste Anmerk.} Diese Vorschriften dienen zugleich zur Bestimmung des
1 \textit{Tisri} oder Neujahrs, welches allezeit 163 Tage nach Ostern des vorhergehenden
Jahres einfällt.

\textit{Zweite Anmerk.} Das Jahr $A$ ist ein gemeines Jahr (von 12 Monaten), wenn
$a<12$, hingegen ein Schaltjahr (von 13 Monaten), wenn $a>11$.

\begin{center}\small
\begin{tabular}{lll}
\textit{Beispiel zu diesen Vorschriften:} & $A=5562$ & $B=1802$\\
& $12A+17=66761$ & $12B+12=21636$\\
\multicolumn{3}{c}{mit 19 dividirt gibt $a=14$}\\
\multicolumn{3}{c}{5562 oder 1802 mit 4 dividirt gibt $b=2$.}\\
Hieraus Werth obiger Formel: & $32.0440932$ & $20.0955877$\\
& $+21.7593852$ & $21.7593852$\\
& $+\phantom{0}0.5$ & $0.5$\\
& $-17.6748903$ & $-\phantom{0}5.7263848$\\\cline{2-3}
& $36.6285881$ & $36.6285881$
\end{tabular}
\end{center}
also $M=36,\quad m=0.6285881$,
\[
\begin{aligned}
M+3A+5b+5&=36+16686+10+5=16737\\
M+3B+5b+1&=36+\phantom{0}5406+10+1=\phantom{0}5453
\end{aligned}
\]
mit 7 dividirt gibt $c=0$.

Da $m$ hier kleiner als $0.89772376$, so kann die Regel III. hier nicht eintreten,
und es ist daher nach IV. Ostern den 36. März alten Styls oder den
48. März neuen Styls, d. i. den 17. April.

In den \textit{meisten} Fällen ist es hinreichend, von obiger Formel nur etwa 2
Decimalstellen zu berechnen.

\begin{center}\rule{.25\linewidth}{.3pt}\end{center}
